% cap1.tex

\chapter{Introducción}
\label{c1} % la etiqueta para referencias

% ============================================================
% Contenido retomado de Proyecto de Título I (PT1), Hito 1 (Hito_1.tex)
% ============================================================

Las infraestructuras financieras alternativas al sistema bancario tradicional
—agrupadas bajo el concepto de \textit{shadow banking} o \textit{non-bank financial
intermediation} (NBFI)— representan hoy USD~238 billones a nivel global, equivalentes
al 49,1\% de los activos financieros mundiales \citeA{fsb_global_2024}. Esta escala
convierte en un problema económico relevante el estudio de cómo las restricciones
regulatorias sobre las redes financieras tradicionales inciden en el bienestar
agregado, y bajo qué condiciones la coexistencia de infraestructuras con
regulaciones distintas —tradicionales y alternativas, incluyendo aquellas basadas en
tecnología \textit{blockchain}— puede recuperar parte o la totalidad de ese bienestar.

Este trabajo se enmarca en la extensión del modelo de redes financieras duales
propuesto por \citeA{venegas_estabilidad_nodate}, quien a su vez extiende el marco
de negociación bilateral en redes de \citeA{jalan_incentive-aware_2024}. El modelo
de Venegas (2024) muestra que, cuando una restricción regulatoria impide el
contrato directo entre dos agentes, la sola competencia entre intermediarios
habilitados recupera solo parcialmente el bienestar perdido, mientras que habilitar
una infraestructura alternativa con regulación distinta —donde el contrato
prohibido en la red tradicional sí puede ejecutarse— logra una recuperación
completa del bienestar agregado. Este resultado, sin embargo, se obtiene bajo
parámetros ilustrativos y no calibrados contra datos reales, lo que limita su
interpretación cuantitativa y su valor como evidencia para política regulatoria.

\textbf{Nota sobre el origen de este capítulo:} la revisión crítica del modelo de
Venegas (2024) que se presenta en el Capítulo~\ref{c2} y las preguntas de
investigación de esta sección fueron desarrolladas originalmente durante
\textit{Proyecto de Título I (PT1, primer semestre de 2026)}, documentadas en
\texttt{Hito\_1.tex} y \texttt{Hito\_2.tex}. Este capítulo las incorpora y actualiza
para reflejar el enfoque acordado con el profesor guía para \textit{Proyecto de
Título II (PT2, segundo semestre de 2026)}, detallado en el
Capítulo~\ref{MCP:subastador}.

\section{Objetivos}
\subsection{Objetivo general}

Evaluar la capacidad del modelo de redes financieras duales de
\citeA{venegas_estabilidad_nodate} para representar y predecir la recuperación de
bienestar observada en un caso de datos reales, aportando así el insumo de
validación empírica que la tesis de Venegas (2024) no desarrolla y que resulta
necesario para su eventual publicación en una revista de finanzas o de
investigación de operaciones.

\subsection{Objetivos específicos}

\begin{enumerate}
    \item Evaluar la factibilidad de calibrar los primitivos del modelo de Venegas
    (2024) —creencias de retorno, matrices de covarianza y coeficientes de aversión
    al riesgo— contra datos reales disponibles a través del Banco Interamericano
    de Desarrollo (BID) y CORFO, en lugar de los valores ilustrativos utilizados
    en el trabajo original.
    \item Implementar y validar numéricamente el modelo base contra los
    \textit{benchmarks} ya publicados por Venegas (2024) (escenarios 3F-WITH/WITHOUT
    y 4F-WITH/WITHOUT), como condición previa a su aplicación empírica.
    \item Aplicar el modelo calibrado al caso de datos reales seleccionado y
    contrastar sus predicciones de bienestar y de reparto de actividad entre redes
    con la evidencia observada.
\end{enumerate}

\section{Preguntas de investigación}

Con base en el marco de Venegas (2024) y bajo el enfoque empírico acordado para
PT2 (Capítulo~\ref{MCP:subastador}), este trabajo busca responder tres preguntas
que el modelo original, calibrado solo con parámetros ilustrativos, no puede
responder por sí mismo:

\begin{enumerate}
    \item ?`Persiste la recuperación (parcial o completa) del bienestar que predice
    el modelo de Venegas (2024) cuando sus parámetros se calibran con datos reales,
    en lugar de valores ilustrativos?
    \item ?`Es consistente el reparto de actividad entre la red tradicional y la
    red alternativa que predice el modelo con el reparto efectivamente observado
    en el caso de datos reales seleccionado?
    \item ?`Cómo se distribuyen, en el caso calibrado, las ganancias de bienestar
    entre los agentes de la red, y coincide esa distribución con la que predice el
    modelo teórico?
\end{enumerate}

Estas preguntas retoman la estructura de las tres preguntas de investigación
formuladas en \texttt{Hito\_2.tex} durante PT1 —que en esa etapa estaban planteadas
sobre una extensión del modelo a agentes heterogéneos en aversión al riesgo y
creencias ($\gamma_i$, $\mu_i$)— pero se reformulan aquí en la versión homogénea
del modelo, con foco en su validación empírica. La extensión a heterogeneidad
queda como línea de trabajo futura, pendiente de decisión con el profesor guía,
y no se aborda en el presente semestre.
\begin{enumerate}
\item Definir objetivo especifico 1
\item obj esp 2
\end{enumerate}


\section{Alcances}

definir lo que se hace y lo que no se va a hacer en la memoria

\section{Metodología}


\begin{enumerate}
\item estudio de la tematica
\item presentación de los modelos econométricos

\end{enumerate}

\newpage
\section{Estructura del documento}

describir la estructura general del documento global



